\documentclass{article}
\usepackage{graphicx} % Required for inserting images
\usepackage{amsmath}
\usepackage{amssymb}
\usepackage{mathtools}
\usepackage{amsthm}
\usepackage{multirow}
\usepackage{natbib}
\title{Physics-Informed Few-step FLow Model}
\author{Khiem Tran}
\date{May 2026}

\begin{document}

\maketitle

\section{Physics-Informed Few-Step Flow for Single/Coupled PDEs}

We propose a generative PDE solver that combines physics-informed transport with few-step flow matching. Let $x_0 \sim p_0$ is the initial state and $x_1 \sim p_1$ be the state of the physics system at time step $\tau$, which we need to predict.

\paragraph{Physics-informed transport path.}
Instead of using the standard linear interpolation between the initial state and the target state, we
learn a physics-aware transport path. Given \(x_0\sim p_0\) and \(x_1\sim p_{1}\), define
\begin{align}
    x_t
    =
    \psi_\phi(t,x_0,x_1)
    =
    (1-t)x_0 + t x_1
    + \alpha(t)\phi(t,x_0,x_1),
    \qquad t\in[0,1],
\end{align}
    
where \(\alpha(0)=\alpha(1)=0\), and \(\alpha(t)=t(1-t)\). This envelope ensures
that the learned correction does not change the endpoints:
\begin{align}
    \psi_\phi(0,x_0,x_1)=x_0,
    \qquad
    \psi_\phi(1,x_0,x_1)=x_1.
\end{align}

The corresponding path velocity is
\begin{align}
    \dot x_t
    =
    \partial_t \psi_\phi(t,x_0,x_1) = (x_1 - x_0) + \dot{\alpha}(t)\phi(t,x_0,x_1) + \alpha(t)\partial_{t} \phi(t, x_0, x_1) .
\end{align}

We introduce the physics residual penalty as $\mathcal{R}(x_t)$, that measure the violation of physics constraints at state $x_t$. Therefore, we use a weighted physics loss
\begin{align}
    \mathcal L_{\rm phy}(\phi)
    =
    \mathbb E_{t,x_0,x_1}
    \left[
        w(t)\cdot\,
        \|\mathcal R(x_t)\|_2^2
    \right],
\end{align}

where \(w(t)= w_0 + t \alpha\) increases toward \(t=1\). 
We further regularize the path-induced velocity field. Let
\begin{align}
     u_\phi(x,t)
    =
    \mathbb E\left[
        \partial_t\psi_\phi(t,x_0,x_1)
        \mid
        \psi_\phi(t,x_0,x_1)=x
    \right]
\end{align}


be the marginal velocity field induced by the path. We penalize its spatial
roughness:
\begin{align}
    \mathcal L_{\rm sm}(\phi)
    =
    \int_0^1
    \mathbb E_{x_t}
    \left[
        \|\nabla_x u_\phi(x_t,t)\|_2^2
    \right]dt.
\end{align}
    

The path is trained by
\begin{align}
     \mathcal L_{\rm path}(\phi)
    =
    \mathcal L_{\rm phy}(\phi)
    +
    \lambda_{\rm sm}\mathcal L_{\rm sm}(\phi),
\end{align}
   

\section{Few-step flow model parameterization}

\paragraph{Few-step flow parameterization.}
After learning the path, we train a few-step flow map model~\citep{boffi2025flow}.  The model predicts the flow map from time \(s\) to \(t\):
\begin{equation}
    X_{s,t}(x_s) = x_t = x_s + (t-s) v_{\theta}(x_s, s, t)
\end{equation}
When \(s = t\), the shortcut $v_{\theta}(x_s, s,s)$ recovers the instantaneous marginal velocity,
\begin{equation}
 v_{\theta}(x_s, s,s) = \frac{d}{ds}(x_s)    
\end{equation}
The flow matching loss is therefore
\begin{equation}
    \mathcal L_{\rm flow}(\theta)
    =
    \mathbb E_{x_t, t \sim U(0,1)}
    \left[
        \left\|
        v_\theta(x_t,t,t)
        -
        \partial_t\psi_\phi(t,x_0,x_1)
        \right\|_2^2
    \right].
\end{equation}
where the path-encoding model $\phi$ is frozen during the flow training process. 
\paragraph{Self-consistency loss.}
To enable one-step and few-step generation, we require the shortcut model to be
self-consistent. We have 
\begin{equation}
    \partial_{t} X_{s,t} (x_s) = \frac{d}{dt}(x_t) = v_{\theta}(x_t, t,t)
\end{equation}
In other hand, we have 
\begin{align}
     \partial_{t} X_{s,t} (x_s) &= \partial_{t} \big( x_s + (t -s)v_\theta(x_s,s,t) \big) \\
     &= v_{\theta}(x_s, s,t) + \partial_t v_{\theta}(x_s, s,t)
\end{align}

Therefor the training objective becomes 
\begin{align}
    \mathcal{L}_{\text{consist}}(\theta) = \mathbb{E}_{s,t} \Big[ \Big\| v_{\theta}(x_s, s,t) + \partial_t v_{\theta}(x_s, s,t) - \text{sg}(v_{\theta}(x_t, t,t))\Big\|^2_2 \Big]
\end{align}
where $\text{sg}$ is the stop gradient operator. 
\bibliography{references}
\bibliographystyle{plainnat}
\end{document}
